\documentclass[12pt,a4paper]{report}
\usepackage{fontspec}
\usepackage{polyglossia}
\setmainlanguage{vietnamese}
\setmainfont{Arial}
\usepackage[a4paper,margin=2.3cm]{geometry}
\usepackage{setspace}
\usepackage{longtable}
\usepackage{booktabs}
\usepackage{hyperref}
\usepackage{xcolor}
\usepackage{listings}
\usepackage{fancyhdr}
\usepackage{enumitem}
\onehalfspacing
\hypersetup{colorlinks=true,linkcolor=blue,urlcolor=blue}
\pagestyle{fancy}
\fancyhf{}
\fancyhead[L]{Quản lý kho hàng -- Nhóm 11}
\fancyhead[R]{Báo cáo đồ án}
\fancyfoot[C]{\thepage}
\lstset{basicstyle=\ttfamily\footnotesize,breaklines=true,frame=single,numbers=left,numberstyle=\tiny,showstringspaces=false,backgroundcolor=\color{gray!6}}
\newcommand{\code}[1]{\texttt{#1}}
\begin{document}
\begin{titlepage}
\centering
{\Large TRƯỜNG ĐẠI HỌC SƯ PHẠM TP. HỒ CHÍ MINH\par}
\vspace{0.4cm}{\large KHOA CÔNG NGHỆ THÔNG TIN\par}\vfill
{\Huge\bfseries BÁO CÁO ĐỒ ÁN CUỐI KỲ\par}
\vspace{0.5cm}{\LARGE\bfseries HỆ THỐNG QUẢN LÝ KHO HÀNG\par}
\vspace{0.4cm}{\Large Môn: Lập trình trên Windows\par}\vfill
\begin{tabular}{ll}
Giảng viên: & ThS. Lê Thanh Thoại\\
Nhóm: & Nhóm 11\\
Năm học: & 2026
\end{tabular}\vfill
TP. Hồ Chí Minh, tháng 11 năm 2026
\end{titlepage}
\pagenumbering{roman}
\chapter*{Lời cam đoan}
Nhóm cam đoan báo cáo mô tả đúng sản phẩm học tập đã xây dựng, bao gồm mã nguồn,
cơ sở dữ liệu, dữ liệu mẫu, test case và hướng dẫn bàn giao.
\chapter*{Lời cảm ơn}
Nhóm cảm ơn giảng viên đã hướng dẫn WinForms, ADO.NET, SQL Server, MVP, kiểm thử
và triển khai ứng dụng Windows.
\chapter*{Tóm tắt}
Đề tài xây dựng ứng dụng desktop quản lý hàng hóa, loại hàng, nhà cung cấp,
phiếu nhập/xuất, tồn kho, lịch sử, CSV, biểu đồ GDI+ và in phiếu. Ứng dụng dùng
MVP kết hợp Repository/Service, SQL Server và transaction.
\tableofcontents
\listoftables
\clearpage\pagenumbering{arabic}

\chapter{Giới thiệu đề tài}
\section{Bối cảnh và nhu cầu}
Kho hàng có nhiều biến động mỗi ngày. Ghi chép thủ công gây sai lệch số lượng,
khó truy vết và chậm lập báo cáo. Hệ thống này tập trung dữ liệu và cung cấp
luồng thao tác thống nhất cho quản lý và nhân viên kho.
\section{Mục tiêu}
\begin{enumerate}
\item Quản lý tập trung hàng hóa, loại hàng và nhà cung cấp.
\item Lưu phiếu nhập/xuất đồng bộ với chi tiết và tồn kho.
\item Cung cấp báo cáo tồn, lịch sử, biểu đồ, CSV và bản in.
\item Bàn giao được source, database, test case, báo cáo và bản Release.
\end{enumerate}
\section{Phạm vi và đối tượng}
Phạm vi gồm một kho nội bộ, hai vai trò người dùng và các nghiệp vụ cốt lõi.
Chưa bao gồm đa kho, mã vạch, hóa đơn điện tử và bán hàng trực tuyến.

\chapter{Phân tích yêu cầu}
\section{Yêu cầu chức năng}
\begin{longtable}{p{0.12\textwidth}p{0.27\textwidth}p{0.48\textwidth}}
\toprule Mã & Nhóm & Mô tả\\\midrule
FR01 & Đăng nhập & Xác thực tài khoản đang hoạt động.\\
FR02 & Phân quyền & Phân biệt Quản lý và Nhân viên kho.\\
FR03 & Loại hàng & Thêm, sửa, xóa, tìm kiếm, kiểm tra ràng buộc.\\
FR04 & Hàng hóa & Quản lý giá, tồn, mức tồn tối thiểu và loại hàng.\\
FR05 & Nhà cung cấp & Quản lý thông tin và tìm kiếm nhà cung cấp.\\
FR06 & Phiếu nhập & Tăng tồn cùng header và chi tiết bằng transaction.\\
FR07 & Phiếu xuất & Kiểm tra đủ tồn trước khi giảm kho.\\
FR08 & Tồn kho & Lọc trạng thái và tính giá trị tồn qua stored procedure.\\
FR09 & Lịch sử & Lọc ngày, loại phiếu, từ khóa và xem chi tiết.\\
FR10 & CSV & Xuất dữ liệu với UTF-8 BOM.\\
FR11 & Biểu đồ & Vẽ cột nhóm hoặc đường cho nhập, xuất, tồn.\\
FR12 & In ấn & PrintPreview và PrintDocument cho phiếu kho.\\
\bottomrule
\end{longtable}
\section{Yêu cầu phi chức năng}
Giao diện phải dễ dùng; lỗi phải hiển thị rõ; truy vấn dùng tham số; dữ liệu được
bảo vệ bằng khóa ngoại, CHECK và transaction; mã nguồn phải phân lớp; chương trình
không phụ thuộc đường dẫn cá nhân.
\section{Người dùng}
Quản lý sử dụng đầy đủ danh mục, nghiệp vụ và báo cáo. Nhân viên kho thực hiện
nhập/xuất và các thao tác được cấp quyền. Tài khoản mẫu là \code{admin/123456}
và \code{nvkho/123456}.

\chapter{Kiến trúc MVP}
\section{Nguyên tắc}
View chỉ phụ trách hiển thị và nhận thao tác; Presenter điều phối; Model mô tả
dữ liệu; Repository truy vấn; Service xử lý CSV, biểu đồ và in. Form giao tiếp
với Presenter qua các interface View.
\section{Cấu trúc thư mục}
\begin{verbatim}
Models/ Views/ Presenters/ Repositories/ Services/
Database.cs  PhienDangNhap.cs  FrmLogin.cs  FrmMain.cs
\end{verbatim}
\section{Phân công}
\begin{longtable}{p{0.2\textwidth}p{0.68\textwidth}}
\toprule Thành viên & Phạm vi\\\midrule
Thành viên 1 & CSDL, Database.cs, đăng nhập, phân quyền, FrmMain, hàng hóa, loại hàng.\\
Thành viên 2 & Nhà cung cấp, phiếu nhập/xuất, SqlTransaction.\\
Thành viên 3 & Tồn kho, lịch sử, CSV, GDI+ và PrintDocument.\\
\bottomrule
\end{longtable}
\section{Luồng xử lý}
Form phát sự kiện; Presenter đọc interface; Repository hoặc Service xử lý;
Presenter yêu cầu View cập nhật. Cách tách này giúp kiểm thử logic mà không
phụ thuộc trực tiếp vào control.

\chapter{Thiết kế cơ sở dữ liệu}
\section{Các bảng}
\begin{longtable}{p{0.25\textwidth}p{0.25\textwidth}p{0.38\textwidth}}
\toprule Bảng & Khóa & Ý nghĩa\\\midrule
NguoiDung & MaNguoiDung & Tài khoản, họ tên, vai trò, trạng thái.\\
LoaiHang & MaLoai & Danh mục loại hàng.\\
NhaCungCap & MaNCC & Thông tin nhà cung cấp.\\
HangHoa & MaHang, MaLoai & Hàng hóa và tồn hiện tại.\\
PhieuKho & MaPhieu & Header nhập/xuất.\\
ChiTietPhieu & MaPhieu, MaHang & Dòng chi tiết, số lượng, đơn giá.\\
\bottomrule
\end{longtable}
\section{Quan hệ}
NguoiDung có nhiều PhieuKho; NhaCungCap có nhiều PhieuKho; PhieuKho có nhiều
ChiTietPhieu; HangHoa có nhiều ChiTietPhieu; LoaiHang có nhiều HangHoa.
\section{Ràng buộc}
Khóa chính chống trùng; khóa ngoại bảo vệ liên kết; CHECK giới hạn vai trò,
loại phiếu, số lượng và đơn giá. Khi xóa loại hàng còn hàng, Presenter từ chối.
\section{Stored procedure}
Hệ thống có các thủ tục lấy danh sách hàng, thống kê nhập xuất tồn, tính tồn kho
chi tiết, tra cứu lịch sử, lấy chi tiết phiếu và lấy header in phiếu.

\chapter{Thiết kế giao diện}
\section{Đăng nhập}
Form có tên đăng nhập, mật khẩu, hiện mật khẩu, đăng nhập và thoát. ErrorProvider
hiển thị trường rỗng; lỗi kết nối nêu rõ cần kiểm tra SQL Server/App.config.
\section{Màn hình chính}
FrmMain là MDI container, có menu hệ thống, danh mục, nghiệp vụ kho, báo cáo và
cửa sổ. StatusStrip hiển thị người dùng, vai trò và thời gian.
\section{Danh mục và phiếu}
Các Form CRUD có vùng nhập, nút thao tác, tìm kiếm và DataGridView. Form phiếu
cho phép thêm nhiều dòng, tính tổng tiền, kiểm tra tồn khi xuất và lưu transaction.
\section{Báo cáo}
FrmThongKeLichSu có tab lịch sử/in phiếu, tồn kho/cảnh báo và biểu đồ. Các chức
năng CSV, PrintPreview, lưu ảnh và lọc dữ liệu được đặt trong Presenter/Service.

\chapter{Cài đặt và triển khai}
\section{Môi trường}
Visual Studio 2022 trở lên, .NET Framework 4.7.2, C\# 7.3 và SQL Server 2019
hoặc SQL Server Express.
\section{Cài database}
Chạy \code{Scripts/QuanLyKhoHangDB.sql} một lần trong SSMS. Với SQLCMD dùng
\code{-f 65001} vì script chứa tiếng Việt UTF-8. Backup trong Database là phương
án thay thế, không chạy đồng thời với script.
\section{Build}
Mở \code{QuanLyKhoHang.slnx}, chọn Release, Build Solution. File chạy nằm trong
\code{SourceCode/QuanLyKhoHang/QuanLyKhoHang/bin/Release/}.
\section{Đóng gói}
Gói nộp gồm SourceCode, Database, Report, TestCases và README. Connection string
nằm trong App.config để thay instance SQL Server trên máy người chấm.

\chapter{Mô tả chức năng}
\section{Đăng nhập và phân quyền}
LoginPresenter kiểm tra đầu vào, gọi TaiKhoanRepository và tạo PhienDangNhap.
FrmMain kiểm soát menu báo cáo theo vai trò.
\section{Hàng hóa và loại hàng}
Presenter kiểm tra rỗng, trùng mã, dữ liệu số và gọi Repository. Lưới được nạp
lại sau mỗi thao tác. Tồn thấp được nhận diện theo mức tối thiểu.
\section{Nhà cung cấp}
Presenter điều phối chế độ thêm/sửa, tìm kiếm, lưu và xóa; dữ liệu được truy vấn
bằng tham số.
\section{Phiếu nhập/xuất}
Phiếu gồm header và nhiều chi tiết. Với xuất kho, Presenter kiểm tra tồn trước
khi lưu và kiểm tra lại trong transaction để tránh biến động đồng thời.
\section{Tồn kho và lịch sử}
Tồn kho gồm giá trị tồn, tổng nhập, tổng xuất và trạng thái. Lịch sử dùng khoảng
ngày nửa mở để bao phủ toàn bộ ngày kết thúc.
\section{CSV, biểu đồ và in}
CSV dùng UTF-8 BOM và escape dữ liệu. GDI+ vẽ cột nhóm/đường có chú giải.
PrintDocument lấy header và chi tiết từ stored procedure để xem trước/in.

\chapter{Kiểm thử}
\section{Chiến lược}
Kiểm thử gồm giao diện, CRUD, nghiệp vụ, CSDL và bàn giao. Mỗi test case có ID,
điều kiện, dữ liệu, thao tác, kết quả mong đợi, thực tế và trạng thái.
\section{Ma trận test case}
\begin{longtable}{p{0.08\textwidth}p{0.2\textwidth}p{0.38\textwidth}p{0.24\textwidth}}
\toprule ID & Chức năng & Dữ liệu/thao tác & Kết quả mong đợi\\\midrule
TC01 & Đăng nhập & admin/123456 & Mở màn hình chính.\\
TC02 & Đăng nhập sai & Mật khẩu sai & Báo lỗi.\\
TC03 & Phân quyền & nvkho mở báo cáo & Giới hạn đúng vai trò.\\
TC04 & Thêm hàng & Mã mới hợp lệ & Lưu và cập nhật lưới.\\
TC05 & Trùng mã & Dùng HH01 lần hai & Từ chối và báo lỗi.\\
TC06 & Xóa loại & Xóa LH01 còn hàng & Từ chối do ràng buộc.\\
TC07 & Nhập kho & Lưu phiếu có chi tiết & Tăng tồn đồng bộ.\\
TC08 & Xuất vượt tồn & Xuất quá HH02 & Không lưu.\\
TC09 & Rollback & Gây lỗi trong transaction & Dữ liệu không đổi.\\
TC10 & Lọc lịch sử & 01/10--20/10/2026 & Đúng khoảng ngày.\\
TC11 & Tồn kho & Lọc cảnh báo/hết hàng & Đúng trạng thái.\\
TC12 & CSV & Xuất DataGridView & Excel đọc đúng tiếng Việt.\\
TC13 & Biểu đồ & Chọn cột/đường & Vẽ đủ số liệu.\\
TC14 & In phiếu & PrintPreview & Đúng header/chi tiết.\\
\bottomrule
\end{longtable}
\section{Kết quả môi trường}
Database thực tế có 2 tài khoản, 5 mặt hàng và 5 phiếu mẫu. Đăng nhập hai vai trò,
các stored procedure tồn kho, lịch sử, thống kê và chi tiết phiếu đã được thực thi.
Bản Release build thành công và khởi động được.

\chapter{Đánh giá và hướng phát triển}
\section{Đáp ứng rubric}
Chức năng bao phủ nghiệp vụ kho; mã nguồn có MVP, Repository và Service; database
có script, backup và dữ liệu; báo cáo có mô tả yêu cầu, kiến trúc, kiểm thử và
phụ lục. Nhóm cần chụp ảnh giao diện và điền kết quả thực tế trước khi nộp.
\section{Hạn chế}
Mật khẩu mẫu đang lưu dạng văn bản; chưa có audit log, đa kho, barcode, backup
tự động và bộ test tự động.
\section{Hướng phát triển}
Băm mật khẩu, logging, nhiều kho, barcode, xuất Excel, phân trang, cấu hình
server bằng giao diện, CI build và kiểm thử tích hợp.

\chapter{Kịch bản demo 7 phút}
\begin{enumerate}
\item Giới thiệu bài toán và người dùng.
\item Đăng nhập tài khoản demo.
\item Demo thêm/sửa/xóa/tìm kiếm.
\item Demo phiếu nhập và phiếu xuất.
\item Demo tồn kho, lịch sử, biểu đồ, CSV hoặc in.
\item Chỉ ra database và kiến trúc MVP.
\item Nêu hạn chế và hướng phát triển.
\end{enumerate}

\appendix
\chapter{Phụ lục mã nguồn}
\lstinputlisting[language={[Sharp]C}]{../SourceCode/QuanLyKhoHang/QuanLyKhoHang/Database.cs}
\lstinputlisting[language={[Sharp]C}]{../SourceCode/QuanLyKhoHang/QuanLyKhoHang/FrmMain.cs}
\lstinputlisting[language={[Sharp]C}]{../SourceCode/QuanLyKhoHang/QuanLyKhoHang/FrmLogin.cs}
\lstinputlisting[language={[Sharp]C}]{../SourceCode/QuanLyKhoHang/QuanLyKhoHang/FrmHangHoa.cs}
\lstinputlisting[language={[Sharp]C}]{../SourceCode/QuanLyKhoHang/QuanLyKhoHang/FrmLoaiHang.cs}
\lstinputlisting[language={[Sharp]C}]{../SourceCode/QuanLyKhoHang/QuanLyKhoHang/PhienDangNhap.cs}
\lstinputlisting[language={[Sharp]C}]{../SourceCode/QuanLyKhoHang/QuanLyKhoHang/Presenters/HangHoaPresenter.cs}
\lstinputlisting[language={[Sharp]C}]{../SourceCode/QuanLyKhoHang/QuanLyKhoHang/Presenters/LoaiHangPresenter.cs}
\lstinputlisting[language={[Sharp]C}]{../SourceCode/QuanLyKhoHang/QuanLyKhoHang/Presenters/LoginPresenter.cs}
\lstinputlisting[language={[Sharp]C}]{../SourceCode/QuanLyKhoHang/QuanLyKhoHang/Presenters/NhaCungCapPresenter.cs}
\lstinputlisting[language={[Sharp]C}]{../SourceCode/QuanLyKhoHang/QuanLyKhoHang/Presenters/PhieuKhoPresenter.cs}
\lstinputlisting[language={[Sharp]C}]{../SourceCode/QuanLyKhoHang/QuanLyKhoHang/Presenters/ThongKeLichSuPresenter.cs}
\lstinputlisting[language={[Sharp]C}]{../SourceCode/QuanLyKhoHang/QuanLyKhoHang/Repositories/HangHoaRepository.cs}
\lstinputlisting[language={[Sharp]C}]{../SourceCode/QuanLyKhoHang/QuanLyKhoHang/Repositories/LoaiHangRepository.cs}
\lstinputlisting[language={[Sharp]C}]{../SourceCode/QuanLyKhoHang/QuanLyKhoHang/Repositories/TaiKhoanRepository.cs}
\lstinputlisting[language={[Sharp]C}]{../SourceCode/QuanLyKhoHang/QuanLyKhoHang/Services/BieuDoGdiHelper.cs}
\lstinputlisting[language={[Sharp]C}]{../SourceCode/QuanLyKhoHang/QuanLyKhoHang/Services/CsvHelperCustom.cs}
\lstinputlisting[language={[Sharp]C}]{../SourceCode/QuanLyKhoHang/QuanLyKhoHang/Services/PhieuKhoPrintHelper.cs}
\lstinputlisting[language={[Sharp]C}]{../SourceCode/QuanLyKhoHang/QuanLyKhoHang/Views/FrmNhaCungCap.cs}
\lstinputlisting[language={[Sharp]C}]{../SourceCode/QuanLyKhoHang/QuanLyKhoHang/Views/FrmPhieuKho.cs}
\lstinputlisting[language={[Sharp]C}]{../SourceCode/QuanLyKhoHang/QuanLyKhoHang/Views/FrmThongKeLichSu.cs}
\lstinputlisting[language={[Sharp]C}]{../SourceCode/QuanLyKhoHang/QuanLyKhoHang/Views/IHangHoaView.cs}
\lstinputlisting[language={[Sharp]C}]{../SourceCode/QuanLyKhoHang/QuanLyKhoHang/Views/ILoaiHangView.cs}
\lstinputlisting[language={[Sharp]C}]{../SourceCode/QuanLyKhoHang/QuanLyKhoHang/Views/ILoginView.cs}
\lstinputlisting[language={[Sharp]C}]{../SourceCode/QuanLyKhoHang/QuanLyKhoHang/Views/INhaCungCapView.cs}
\lstinputlisting[language={[Sharp]C}]{../SourceCode/QuanLyKhoHang/QuanLyKhoHang/Views/IPhieuKhoView.cs}
\lstinputlisting[language={[Sharp]C}]{../SourceCode/QuanLyKhoHang/QuanLyKhoHang/Views/IThongKeLichSuView.cs}
\chapter{Phụ lục SQL}
\lstinputlisting[language=SQL]{../Scripts/QuanLyKhoHangDB.sql}
\lstinputlisting[language=SQL]{../Database/01_CreateDatabase.sql}
\lstinputlisting[language=SQL]{../Database/03_ThanhVien3_StoredProcedures_DuLieu.sql}
\chapter{Checklist bàn giao}
\begin{longtable}{p{0.08\textwidth}p{0.65\textwidth}p{0.15\textwidth}}
\toprule STT & Hạng mục & Xác nhận\\\midrule
1 & Source mở được solution và build Release & \\
2 & Database script, backup và dữ liệu mẫu & \\
3 & Tài khoản demo và phân quyền & \\
4 & CRUD danh mục và ràng buộc & \\
5 & Phiếu nhập/xuất transaction & \\
6 & Tồn kho, lịch sử, CSV, biểu đồ, in & \\
7 & Test case và ảnh minh chứng & \\
8 & Báo cáo đủ bố cục và phụ lục & \\
\bottomrule
\end{longtable}
\end{document}
